% =====================================================================================
%  preamble.tex  —  design system for the "Mathematics for AI Safety" lecture notes
%  Engine: XeLaTeX (or LuaLaTeX).  Build with:  latexmk -xelatex main.tex
% =====================================================================================

% ---------- fonts ---------------------------------------------------------------------
\usepackage{fontspec}
\usepackage{unicode-math}
\setmainfont{Libertinus Serif}[Ligatures=TeX, Numbers=Lining]
\setsansfont{FiraSans}[Extension=.otf, UprightFont=*-Regular, BoldFont=*-SemiBold, ItalicFont=*-Italic,
  BoldItalicFont=*-SemiBoldItalic, Scale=MatchLowercase, Ligatures=TeX, Numbers=Lining]
\setmonofont{FiraMono}[Extension=.otf, UprightFont=*-Regular, BoldFont=*-Medium, Scale=MatchLowercase]
\setmathfont{Libertinus Math}
\newfontfamily\headsf{FiraSans}[Extension=.otf, UprightFont=*-Medium, BoldFont=*-SemiBold, ItalicFont=*-MediumItalic,
  Scale=MatchLowercase, Ligatures=TeX, Numbers=Lining]
\newfontfamily\lightsf{FiraSans}[Extension=.otf, UprightFont=*-Light, BoldFont=*-Regular, ItalicFont=*-LightItalic,
  Scale=MatchLowercase, Ligatures=TeX, Numbers=Lining]

% ---------- page geometry -------------------------------------------------------------
\usepackage[a4paper, left=22mm, textwidth=128mm, marginparsep=7mm, marginparwidth=42mm,
            top=27mm, bottom=27mm, headheight=14pt, headsep=10mm, footskip=12mm]{geometry}
\usepackage{microtype}
\usepackage{setspace}
\setstretch{1.075}
\usepackage[english]{babel}
\usepackage{csquotes}
\setlength{\parindent}{0pt}
\setlength{\parskip}{0.62\baselineskip plus 1pt minus 1pt}

% ---------- colours -------------------------------------------------------------------
\usepackage[dvipsnames,svgnames,table]{xcolor}
\definecolor{ink}{HTML}{14213D}        % headings, strong text
\definecolor{accent}{HTML}{2B59C3}     % definitions, links
\definecolor{teal}{HTML}{0E7C7B}       % theorems / propositions
\definecolor{amber}{HTML}{B8730B}      % intuition / key ideas
\definecolor{crimson}{HTML}{B3261E}    % cautions
\definecolor{plum}{HTML}{6C3FA0}       % open questions
\definecolor{slate}{HTML}{475569}      % remarks, margin notes
\definecolor{moss}{HTML}{3F7D3A}       % examples
\definecolor{rule}{HTML}{C9D1DE}       % hairlines

% ---------- maths ---------------------------------------------------------------------
\usepackage{mathtools}
\usepackage{amsthm}
\usepackage{bm}
\usepackage{enumitem}
\usepackage{booktabs}
\usepackage{array}
\usepackage{multirow}
\usepackage{longtable}
\usepackage{xspace}
\usepackage{xparse}

\DeclareMathOperator*{\argmin}{arg\,min}
\DeclareMathOperator*{\argmax}{arg\,max}
\DeclareMathOperator{\Var}{Var}
\DeclareMathOperator{\Cov}{Cov}
\DeclareMathOperator{\supp}{supp}
\DeclareMathOperator{\KL}{KL}
\DeclareMathOperator{\RLCT}{RLCT}
\DeclareMathOperator{\ran}{ran}
\DeclareMathOperator{\diag}{diag}
\newcommand{\R}{\mathbb{R}}
\newcommand{\N}{\mathbb{N}}
\newcommand{\E}{\mathbb{E}}
\renewcommand{\Pr}{\mathbb{P}}
\newcommand{\dd}{\mathrm{d}}
\newcommand{\eps}{\varepsilon}
\newcommand{\ip}[2]{\left\langle #1,\,#2\right\rangle}
\newcommand{\norm}[1]{\left\lVert #1\right\rVert}
\newcommand{\abs}[1]{\left\lvert #1\right\rvert}
\newcommand{\Hil}{L^2(\mu)}
\newcommand{\NTK}{\mathrm{NTK}}
\newcommand{\simeqq}{\asymp}
\newcommand{\ind}{\mathbf{1}}
\newcommand{\bits}{\text{bits}}
\newcommand{\Ocal}{\mathcal{O}}
\newcommand{\Kc}{\mathsf{K}}                 % Kolmogorov complexity
\newcommand{\rlct}{\lambda_{\mathrm{RLCT}}}  % real log canonical threshold
\newcommand{\Pbayes}{P_{\mathrm{B}}}
\newcommand{\Hent}{\mathrm{H}}
\newcommand{\sig}{\sigma}
\DeclareMathOperator{\softmax}{softmax}
\newcommand{\vect}[1]{\bm{#1}}
\renewcommand{\qedsymbol}{\textcolor{accent}{\raisebox{0.02ex}{\rule{0.58em}{0.58em}}}}
\usepackage{cancel}

% ---------- graphics, TikZ, pgfplots --------------------------------------------------
\usepackage{graphicx}
\usepackage{tikz}
\usetikzlibrary{arrows.meta, calc, positioning, shapes.geometric, shapes.misc, fit, backgrounds, decorations.pathreplacing,
                decorations.markings, patterns, shadows.blur, shapes.symbols, intersections, through, bending, matrix, mindmap}
\usepackage{pgfplots}
\pgfplotsset{compat=1.18}
\usepgfplotslibrary{fillbetween, groupplots, colorbrewer}
\pgfplotsset{
  notes/.style={
    width=\linewidth, height=0.64\linewidth, scale only axis=false,
    axis lines=left, axis line style={-{Latex[length=2mm,width=1.4mm]}, ink!70, line width=0.5pt},
    tick style={ink!70, line width=0.4pt},
    every tick label/.append style={font=\sffamily\scriptsize, text=ink!85},
    label style={font=\sffamily\small, text=ink},
    legend style={font=\sffamily\scriptsize, draw=none, fill=white, fill opacity=0.85, text opacity=1,
                  inner sep=2pt, row sep=1pt, rounded corners=1pt},
    grid=major, major grid style={rule!70, line width=0.3pt},
    minor grid style={rule!40, line width=0.2pt},
    every axis plot/.append style={line width=1.1pt},
    title style={font=\sffamily\small\bfseries, text=ink},
    unbounded coords=jump,
    table/col sep=space,
  }
}
\tikzset{
  every picture/.append style={>={Latex[length=2mm,width=1.4mm]}},
  smallnode/.style={font=\sffamily\footnotesize, text=ink, align=center, inner sep=3pt},
  card/.style={rounded corners=2pt, draw=#1, fill=#1!7, line width=0.7pt, font=\sffamily\footnotesize, text=ink,
               align=center, inner sep=4pt},
  card/.default=accent,
}

% ---------- links & references --------------------------------------------------------
\usepackage[authoryear, round]{natbib}
\setlength{\bibsep}{3pt}
\usepackage[colorlinks=true, linkcolor=accent, citecolor=teal, urlcolor=accent, breaklinks=true,
            pdfauthor={Unofficial notes on lectures by Y. Liokumovich (MAT1510, Fields Institute / University of Toronto)}, pdfsubject={Lecture notes, Lectures 1--8: capability trajectories, scaling laws, generalization and compression, recursive self-improvement, Bayesian learning and grokking, LLM training, emergent misalignment, jailbreaks}, pdfkeywords={AI safety, scaling laws, neural tangent kernel, generalization, Solomonoff induction, singular learning theory, software intelligence explosion, Bayesian learning, grokking, language models, RLHF, emergent misalignment, jailbreaks},
            pdftitle={Mathematics for AI Safety: Lecture Notes, Lectures 1--8},
            bookmarksnumbered=true]{hyperref}
\usepackage{url}
\urlstyle{same}
\Urlmuskip=0mu plus 1mu

% ---------- margin notes (sidenotes + clickable YouTube timestamps) --------------------
\usepackage{marginnote}
\renewcommand*{\raggedleftmarginnote}{\raggedright}
\renewcommand*{\marginfont}{\sffamily\scriptsize\color{slate}\setstretch{1.02}}
\reversemarginpar   % not used: single-sided, margin on the right
\normalmarginpar

% current lecture video id (set at the start of every lecture chapter)
\newcommand{\lecturevideo}{}
\newcommand{\setlecturevideo}[1]{\renewcommand{\lecturevideo}{#1}}

% \ts{mm:ss} or \ts{h:mm:ss}  ->  link to that moment of the current lecture video
\newcommand{\tsfinish}[4]{%
  \href{https://www.youtube.com/watch?v=\lecturevideo&t=\the\numexpr(#1)*3600+(#2)*60+(#3)\relax s}{#4}}
\def\tsparse#1:#2:#3\relax{%
  \ifx\relax#3\relax\tsfinish{0}{#1}{#2}{#1:#2}\else\tsparseB{#1}{#2}#3\relax\fi}
\def\tsparseB#1#2#3:\relax{\tsfinish{#1}{#2}{#3}{#1:#2:#3}}
\newcommand{\tsraw}[1]{\tsparse#1:\relax}
\newcommand{\ts}[1]{\hbox{\sffamily\scriptsize\color{accent}$\blacktriangleright$\,\tsraw{#1}}}
% a margin timestamp aligned with the current line
\newcommand{\tsm}[1]{\marginnote{\ts{#1}}[-0.15\baselineskip]}
% a general margin remark
\newcommand{\mnote}[1]{\marginnote{#1}[-0.15\baselineskip]}

% ---------- boxes ---------------------------------------------------------------------
\usepackage[most]{tcolorbox}
\tcbset{
  notebox/.style n args={2}{% colour, tint
    enhanced jigsaw, breakable, sharp corners,
    colback=#1!#2!white, colframe=#1, boxrule=0pt, frame hidden,
    borderline west={2.6pt}{0pt}{#1},
    left=9pt, right=8pt, top=5pt, bottom=5pt, before skip=9pt, after skip=9pt,
    fontupper=\normalfont, parskip=0.5\baselineskip,
    overlay unbroken and first={},
  },
}
% run-in box head:  \boxhead{colour}{Label}{number or empty}{optional title or -NoValue-}
\NewDocumentCommand{\boxhead}{m m m m}{%
  {\headsf\bfseries\color{#1}#2\IfBlankTF{#3}{}{\,#3}}%
  \IfNoValueTF{#4}{}{\ \textit{(#4)}}{\color{#1}\bfseries.}\enspace}

% counter shared by all numbered boxes
\newcounter{box}[chapter]
\renewcommand{\thebox}{\thechapter.\arabic{box}}

\NewDocumentEnvironment{defn}{o}{%
  \refstepcounter{box}\begin{tcolorbox}[notebox={accent}{6}]\boxhead{accent}{Definition}{\thebox}{#1}}{\end{tcolorbox}}
\NewDocumentEnvironment{thm}{o}{%
  \refstepcounter{box}\begin{tcolorbox}[notebox={teal}{7}]\boxhead{teal}{Theorem}{\thebox}{#1}\itshape}{\end{tcolorbox}}
\NewDocumentEnvironment{prop}{o}{%
  \refstepcounter{box}\begin{tcolorbox}[notebox={teal}{7}]\boxhead{teal}{Proposition}{\thebox}{#1}\itshape}{\end{tcolorbox}}
\NewDocumentEnvironment{lem}{o}{%
  \refstepcounter{box}\begin{tcolorbox}[notebox={teal}{7}]\boxhead{teal}{Lemma}{\thebox}{#1}\itshape}{\end{tcolorbox}}
\NewDocumentEnvironment{claim}{o}{%
  \refstepcounter{box}\begin{tcolorbox}[notebox={teal}{7}]\boxhead{teal}{Claim}{\thebox}{#1}\itshape}{\end{tcolorbox}}
\NewDocumentEnvironment{exmp}{o}{%
  \refstepcounter{box}\begin{tcolorbox}[notebox={moss}{6}]\boxhead{moss}{Example}{\thebox}{#1}}{\end{tcolorbox}}
\NewDocumentEnvironment{remark}{o}{%
  \begin{tcolorbox}[notebox={slate}{5}]\boxhead{slate}{Remark}{}{#1}}{\end{tcolorbox}}
\NewDocumentEnvironment{intuition}{o}{%
  \begin{tcolorbox}[notebox={amber}{8}]\boxhead{amber}{Intuition}{}{#1}}{\end{tcolorbox}}
\NewDocumentEnvironment{caution}{o}{%
  \begin{tcolorbox}[notebox={crimson}{5}]\boxhead{crimson}{Caution}{}{#1}}{\end{tcolorbox}}
\NewDocumentEnvironment{openq}{o}{%
  \begin{tcolorbox}[notebox={plum}{6}]\boxhead{plum}{Open question}{}{#1}}{\end{tcolorbox}}
\NewDocumentEnvironment{discussion}{o}{%
  \begin{tcolorbox}[notebox={slate}{4}, colback=black!3!white]\boxhead{slate}{From the discussion}{}{#1}}{\end{tcolorbox}}
\NewDocumentEnvironment{beyond}{o}{%
  \begin{tcolorbox}[notebox={amber}{3}, colback=amber!4!white, borderline west={2.6pt}{0pt}{amber!60!white}, borderline west={0.8pt}{3.4pt}{amber!60!white}]%
  \boxhead{amber}{Beyond the lecture}{}{#1}\textsuperscript{\textcolor{amber}{\dag}}\enspace}{\end{tcolorbox}}

% key take-aways box (deep navy, white text)
\NewDocumentEnvironment{takeaways}{O{Key take-aways} }{%
  \begin{tcolorbox}[enhanced jigsaw, breakable, colback=ink, colframe=ink, boxrule=0pt, arc=1.6mm,
    left=12pt, right=12pt, top=9pt, bottom=9pt, before skip=14pt, after skip=10pt, coltext=white,
    fontupper=\normalfont, parskip=0.35\baselineskip,
    title={\headsf\bfseries\large #1}, coltitle=white, colbacktitle=ink, fonttitle=\headsf,
    attach title to upper={\par\smallskip},
    overlay={\fill[white, opacity=0.08] ([xshift=-1.6mm]frame.north east) circle (9mm);}]%
  \setlist[itemize]{leftmargin=1.3em, itemsep=0.2em, topsep=0.2em, label=\textcolor{amber!70!white}{\large$\blacktriangleright$}}%
  \hypersetup{urlcolor=white}\color{white}}{\end{tcolorbox}}

% a quotation box
\NewDocumentEnvironment{quotebox}{m }{%
  \begin{tcolorbox}[enhanced jigsaw, breakable, colback=white, colframe=rule, boxrule=0.6pt, sharp corners,
    left=12pt, right=12pt, top=7pt, bottom=7pt, before skip=9pt, after skip=9pt, fontupper=\itshape,
    overlay={\node[font=\Huge\bfseries, text=rule, anchor=north west, inner sep=0pt] at ([xshift=2pt,yshift=-1pt]frame.north west) {\textquotedblleft};}]%
  \setlength{\leftskip}{6pt}}%
  {\par\hfill{\normalfont\sffamily\scriptsize\color{slate}#1}\end{tcolorbox}}

% proofs
\renewcommand{\proofname}{\headsf\color{slate}Proof}
\newcommand{\proofsketchname}{\headsf\color{slate}Proof idea}
\newenvironment{proofidea}{\par\noindent{\proofsketchname.}\enspace}{\par}

% exercise lists
\NewDocumentEnvironment{exercise}{o}{%
  \begin{tcolorbox}[notebox={ink}{4}, colback=ink!4!white, unbreakable]\boxhead{ink}{Exercise}{}{#1}}{\end{tcolorbox}}

% ---------- lists ---------------------------------------------------------------------
\setlist[itemize]{leftmargin=1.35em, itemsep=0.15em, topsep=0.25em, parsep=0.2em,
                  label=\textcolor{accent}{\raisebox{0.12ex}{\scriptsize$\blacktriangleright$}}}
\setlist[enumerate]{leftmargin=1.6em, itemsep=0.15em, topsep=0.25em, parsep=0.2em, label=\textcolor{accent}{\bfseries\arabic*.}}
\setlist[description]{font=\headsf\color{ink}, leftmargin=0pt, itemsep=0.2em}

% ---------- captions ------------------------------------------------------------------
\usepackage{caption}
\usepackage{float}
\captionsetup{font={small,sf}, labelfont={bf,color=accent}, labelsep=period, justification=raggedright,
              singlelinecheck=false, skip=6pt, margin=0pt}
\setlength{\textfloatsep}{14pt plus 2pt minus 4pt}
\setlength{\intextsep}{12pt plus 2pt minus 2pt}

% ---------- KOMA structure: chapters = lectures ---------------------------------------
\KOMAoptions{headings=big}
\setkomafont{disposition}{\headsf\color{ink}}
\setkomafont{section}{\headsf\Large\color{ink}}
\setkomafont{subsection}{\headsf\large\color{ink}}
\setkomafont{subsubsection}{\headsf\normalsize\color{slate}}
\setkomafont{paragraph}{\headsf\normalsize\color{accent}}
\renewcommand*{\sectionformat}{\textcolor{accent}{\thesection}\enskip}
\renewcommand*{\subsectionformat}{\textcolor{accent}{\thesubsection}\enskip}
\RedeclareSectionCommand[beforeskip=1.9\baselineskip, afterskip=0.7\baselineskip]{section}
\RedeclareSectionCommand[beforeskip=1.4\baselineskip, afterskip=0.45\baselineskip]{subsection}
\RedeclareSectionCommand[beforeskip=1.1\baselineskip, afterskip=0.3\baselineskip]{subsubsection}
\setkomafont{descriptionlabel}{\headsf\color{ink}}

% chapter heading: big numeral + title
\renewcommand*{\chapterformat}{\thechapter}
\renewcommand*{\chapterlinesformat}[3]{%
  \ifstr{#1}{chapter}{\ifstrempty{#2}{\par{\headsf\bfseries\fontsize{25}{28}\selectfont\color{ink}#3}\par}{%
    \par\vspace*{-3mm}%
    \begin{tikzpicture}
      \node[anchor=south west, inner sep=0pt] (num) at (0,0)
        {\fontsize{76}{76}\selectfont\headsf\bfseries\color{accent!85!white}#2};
      \node[anchor=south west, inner sep=0pt] (ttl) at ([xshift=5.5mm,yshift=-0.5mm]num.south east)
        {\parbox[b]{100mm}{\raggedright{\headsf\scshape\small\color{slate}\addfontfeature{LetterSpace=8}Lecture \thechapter}\\[3pt]%
         {\headsf\bfseries\fontsize{25}{28}\selectfont\color{ink}#3}}};
      \draw[accent, line width=1.4pt] ([xshift=-2.6mm,yshift=1mm]ttl.south west) -- ([xshift=-2.6mm,yshift=-1mm]ttl.north west);
    \end{tikzpicture}%
  }}{#3}}
\renewcommand*{\chapterheadstartvskip}{\vspace*{0\baselineskip}}
\renewcommand*{\chapterheadendvskip}{\vspace*{0.8\baselineskip}}

% header and footer
\usepackage[headsepline=0.4pt, footsepline=0pt, plainfootsepline=false]{scrlayer-scrpage}
\addtokomafont{pageheadfoot}{\sffamily\scriptsize\color{slate}}
\addtokomafont{pagenumber}{\sffamily\scriptsize\bfseries\color{ink}}
\setkomafont{headsepline}{\color{rule}}
\pagestyle{scrheadings}
\clearpairofpagestyles
\ihead{MAT1510 \textperiodcentered{} Mathematics for AI Safety}
\ohead{\headmark}
\cfoot*{\pagemark}
\ofoot{}
\automark[section]{chapter}
\renewcommand*{\chaptermarkformat}{Lecture \thechapter\enskip}
\renewcommand*{\sectionmarkformat}{\thesection\enskip}
\pagestyle{scrheadings}
\setkomafont{pageheadfoot}{\sffamily\scriptsize\color{slate}}
\renewcommand*{\chapterpagestyle}{scrplain}
\setkomafont{pagenumber}{\sffamily\scriptsize\bfseries\color{ink}}
\cfoot[\pagemark]{\pagemark}

% table of contents
\newcommand{\tocLecture}[1]{\textcolor{accent}{Lecture\,#1}}
\DeclareTOCStyleEntry[indent=0pt, beforeskip=1.1em plus 2pt, numwidth=5.4em, dynnumwidth=false,
  entryformat=\headsf\bfseries\large\color{ink}, entrynumberformat=\tocLecture,
  pagenumberformat=\headsf\bfseries\color{ink}]{tocline}{chapter}
\DeclareTOCStyleEntry[indent=5.4em, beforeskip=0.1em, numwidth=2.3em, dynnumwidth=false,
  entryformat=\small\normalfont\color{ink!90}]{tocline}{section}
\DeclareTOCStyleEntry[indent=7.7em, beforeskip=0pt, numwidth=3.0em, dynnumwidth=false,
  entryformat=\footnotesize\normalfont\color{ink!70}]{tocline}{subsection}
\setcounter{tocdepth}{2}
\setcounter{secnumdepth}{3}

% unnumbered section with a TOC entry (and a working hyperlink anchor)
\newcommand{\tocsection}[1]{\needspace{16\baselineskip}\section*{#1}\phantomsection\addcontentsline{toc}{section}{\protect\nonumberline #1}}
% ---------- small helpers --------------------------------------------------------------
\newcommand{\term}[1]{\textbf{\color{ink}#1}}               % first use of a term
\newcommand{\say}[1]{\textit{#1}}                            % paraphrase of something said in lecture
\newcommand{\flag}{\textsuperscript{\textcolor{crimson}{\ddag}}}  % ASR-uncertain item
\newcommand{\dagmark}{\textsuperscript{\textcolor{amber}{\dag}}} % added from literature
\newcommand{\vid}[2]{\href{https://www.youtube.com/watch?v=#1}{#2}}
\newcommand{\kbd}[1]{\texttt{#1}}
\newcommand{\sfrac}[2]{\mathord{\raisebox{0.15ex}{\scriptsize$#1$}\!/\!\raisebox{-0.2ex}{\scriptsize$#2$}}}
\newcommand{\lectureinfo}[4]{%
  \begin{tcolorbox}[enhanced jigsaw, colback=ink!4!white, colframe=ink!4!white, boxrule=0pt, arc=1.4mm,
      left=9pt, right=9pt, top=5pt, bottom=5pt, before skip=4pt, after skip=8pt, fontupper=\sffamily\scriptsize]
    \begin{tabularx}{\linewidth}{@{}Xr@{}}
      \textcolor{slate}{\textbf{Date}}\enspace #1 \quad \textcolor{slate}{\textbf{Length}}\enspace #2 &
      \textcolor{accent}{$\blacktriangleright$}\ \href{https://www.youtube.com/watch?v=#3}{Watch the recording}\\
      \multicolumn{2}{@{}>{\raggedright\arraybackslash}X@{}}{\textcolor{slate}{\textbf{Instructor}}\enspace Yevgeny Liokumovich (University of Toronto)\quad \textcolor{slate}{\textbf{Theme}}\enspace #4}
    \end{tabularx}
  \end{tcolorbox}}
\usepackage{tabularx}
\renewcommand{\arraystretch}{1.18}
% lecture header for a guest lecture:  date, length, video id, speaker, theme
\newcommand{\lectureinfoguest}[5]{%
  \begin{tcolorbox}[enhanced jigsaw, colback=ink!4!white, colframe=ink!4!white, boxrule=0pt, arc=1.4mm,
      left=9pt, right=9pt, top=5pt, bottom=5pt, before skip=4pt, after skip=8pt, fontupper=\sffamily\scriptsize]
    \begin{tabularx}{\linewidth}{@{}Xr@{}}
      \textcolor{slate}{\textbf{Date}}\enspace #1 \quad \textcolor{slate}{\textbf{Length}}\enspace #2 &
      \textcolor{accent}{$\blacktriangleright$}\ \href{https://www.youtube.com/watch?v=#3}{Watch the recording}\\
      \multicolumn{2}{@{}>{\raggedright\arraybackslash}X@{}}{\textcolor{slate}{\textbf{Guest speaker}}\enspace #4\quad \textcolor{slate}{\textbf{Theme}}\enspace #5}
    \end{tabularx}
  \end{tcolorbox}}

% an "at a glance" box
\NewDocumentEnvironment{glance}{}{%
  \begin{tcolorbox}[enhanced jigsaw, breakable, colback=white, colframe=accent!60!white, boxrule=0.7pt, arc=1.6mm,
    left=11pt, right=11pt, top=6pt, bottom=6pt, before skip=8pt, after skip=10pt, fontupper=\normalfont,
    title={\headsf\bfseries\small\scshape In this lecture}, coltitle=accent, colbacktitle=white, fonttitle=\headsf,
    attach title to upper={\par\vspace{2pt}}, bottomtitle=0pt, toptitle=0pt]}{\end{tcolorbox}}

% lecture-map: a horizontal bar of sections, proportional to the video duration
\newlength{\mapunit}
% \mapbar{total seconds}{tick step in seconds}{TikZ commands using \mapseg / \maplabel}
\newcommand{\mapbar}[3]{%
  \pgfmathsetlength{\mapunit}{\linewidth/#1}%
  \makebox[\linewidth][c]{\begin{tikzpicture}[x=\mapunit, y=1mm]
    #3
    \draw[slate!55, line width=0.35pt] (0,-0.9) -- (#1,-0.9);
    \foreach \s in {0,#2,...,#1} {%
      \pgfmathtruncatemacro{\mm}{\s/60}%
      \draw[slate!55, line width=0.35pt] (\s,-0.9) -- (\s,-1.9);%
      \node[font=\sffamily\fontsize{6.8}{7.5}\selectfont, text=slate, anchor=north, inner sep=1pt] at (\s,-1.9)
        {\ifnum\mm=0 0\else\ifnum\mm<60 \mm$'$\else\ifnum\mm=60 1h\else 1h\the\numexpr\mm-60\relax\fi\fi\fi};}%
  \end{tikzpicture}}}
% \mapseg{start}{end}{colour}
\newcommand{\mapseg}[3]{\fill[#3!80!white] (#1,0) rectangle (#2,5.2); \draw[white, line width=0.8pt] (#1,0) -- (#1,5.2);}
% \maplabel{x}{text}{height}{anchor}  -- label above the bar with a thin leader
\newcommand{\maplabel}[4]{\draw[slate!70, line width=0.3pt] (#1,5.2) -- (#1,5.2+#3);
  \node[font=\sffamily\fontsize{6.4}{7.2}\selectfont, text=ink, anchor=#4, inner sep=1.2pt, align=left] at (#1,5.2+#3) {#2};}
% label written inside a wide segment
\newcommand{\mapin}[2]{\node[font=\sffamily\fontsize{6.4}{7}\selectfont\bfseries, text=white, inner sep=0pt] at (#1,2.6) {#2};}

\usepackage{etoolbox}
\usepackage{lipsum}
\usepackage{fancyvrb}
\usepackage{fontawesome5}
\usepackage{needspace}

% balanced bibliography heading
\renewcommand{\bibsection}{\chapter*{References}\phantomsection\addcontentsline{toc}{chapter}{\protect\nonumberline References}\markboth{References}{References}}
\setlength{\bibhang}{1.4em}

\emergencystretch=1.5em
\clubpenalty=4500
\widowpenalty=4500
\displaywidowpenalty=3000
\brokenpenalty=2000
